% =====================================================================================
% WHITE PAPER discussion (user 2026-11-12; restructured after a colleague's review: one question, transitions, what it
% means for experimentalists and engineers, "first step towards X"): implications, idealisations, extensions and vision.
% No functional / DFT details.
% =====================================================================================

The results trace one chain from process to device. The cation fixes the tilt pattern; the anneal and the substrate
leave the film under 0.2 to \SI{0.4}{\percent} in-plane tension, within the range measured in films, without changing
that pattern; and the tension acts more on the lifetime of the film than on its band gap. Four points follow for
experiment and manufacturing.

\paragraph{Choice of cation.} A cation with hydrogen-bonding groups at both ends (FA) selects in-phase tilting, one
anchored at a single end (MA) antiphase tilting: a criterion for anticipating the tilt pattern of new cations and
mixtures.

\paragraph{Band gap.} For the untilted hybrids the residual strain of processing shifts the gap by about \SI{10}{meV},
less than its thermal fluctuation at the operating temperature. Shifts of several tens of meV measured on such films
therefore point to causes other than processing strain, such as composition or carrier concentration.

\paragraph{Lifetime.} Residual tension speeds up iodide migration in \CsSnI{}, by 1.2 to 1.5 times at the strain left
by processing and about twice at \SI{1}{\percent}, while the hybrids are hardly affected. For \CsSnI{} films, processing
that lowers the tension should slow this route of degradation: a lower anneal temperature, a substrate or transport
layer whose thermal expansion is closer to that of the film, or a transport layer deposited hot so that it compresses
the film on cooling, as demonstrated for CsPbI$_2$Br.\cite{ref_processing_temp}

\paragraph{Black \CsSnI{}.} At the operating temperature the yellow phase is the stable one, and the residual stress of
processing shifts the balance by only 2 to \SI{4}{meV} per formula unit. Keeping a film black therefore requires
either in-plane strains near \SI{1}{\percent} or protection that keeps the conversion from nucleating.

\paragraph{Idealisations.} The simulations describe an idealised film. Each cell is a periodic single crystal a few
nanometres across, without grains, grain boundaries, free surfaces or interfaces with the transport layers, although
these relieve strain, nucleate tilt domains and offer fast paths for ions. The substrate enters only through a fixed
in-plane stress or a rigid in-plane cell; a real substrate also contracts, and the bond to the film relaxes over hours,
so the rigid case is an upper bound of the residual strain. The films are cooled many orders of magnitude faster than in
processing, and slow, reconstructive changes such as the conversion of black to yellow \CsSnI{} are treated through
free energies rather than followed directly. The chemistry is
restricted to the perovskite and its cation and iodide vacancies: oxygen, moisture, the oxidation of Sn(II) to Sn(IV),
light and excess charge carriers, which drive the degradation of tin films in practice,\cite{ref_sn_instability} are
not included.

\paragraph{Extensions and vision.} These idealisations can be removed step by step. Adding oxygen and nitrogen will
compare processing and operation in air with inert nitrogen atmospheres and give access to the first steps of Sn(II)
oxidation. Polycrystalline cells with grain boundaries, built from cells of tens of thousands of atoms that already
form tilt domains, will show how boundaries relieve the residual strain, pin the tilt variant that forms on cooling
and channel iodide migration. An explicit substrate and transport layer will replace the idealised boundary
conditions by the real interface and its mismatch of thermal expansion. This white paper is the first step toward a
virtual process line for perovskite absorbers. With these extensions the same framework can rank process windows
(anneal temperature, cooling rate, substrate, transport layer and atmosphere) by the band gap and lifetime they
produce, and test compositions before they are synthesised. For manufacturers this replaces part of the trial and
error of process tuning by a screened search; for experimentalists it provides the structural picture needed to
interpret diffraction, photoluminescence and ion-migration measurements on strained films.
